%% Information Systems (Elsevier) -- research article.
%% Compile: python scripts/build_journal.py --paper specification_surfaces
%%
%% ROUND TWENTY-ONE.  Three reviews have been answered: a referee's report for
%% Information Systems (response_to_referee.md), a developmental review of the
%% round-nineteen manuscript (response_to_blueprint.md), and a second
%% developmental review of the round-twenty manuscript, whose twelve major
%% comments this round answers (response_to_review21.md).  REFEREE-LOG.md
%% carries the objection behind every correction.
%%
%% THE APPENDICES ARE A SEPARATE DOCUMENT, paper/supplement.tex, built by the
%% same command.  Cross-references between the two resolve through `xr'.
%%
%% EVERY NUMBER IN THIS DOCUMENT IS A MACRO defined in paper/numbers.tex,
%% which scripts/make_numbers.py generates from results/*.csv.  A number
%% therefore cannot disagree between the abstract, a table and the
%% conclusion: there is one of it.  scripts/verify_paper.py checks each
%% macro against the result file it came from.
\documentclass[preprint,12pt]{elsarticle}

%%  Latin Modern carries the bold and small-caps typewriter shapes that
%%  Computer Modern lacks -- which is what thirty font-substitution warnings
%%  in an earlier build of this round were -- and is scalable, so the PDF
%%  embeds outlines rather than bitmaps.
\usepackage{lmodern}
\usepackage{amsmath}
\usepackage{amssymb}
\usepackage{amsthm}
\usepackage{graphicx}
\usepackage{booktabs}
%%  `array' supplies the >{...} column modifier, which is what lets a narrow
%%  text column in a wide table be set ragged right instead of justified.
\usepackage{array}
\usepackage[hyphens]{url}
\urlstyle{rm}
\usepackage[protrusion=true,expansion=false]{microtype}
%%  The supplement is a separate document.  `xr' resolves a reference to one
%%  of its sections or tables to the number that document settled on, so the
%%  article can point at Section S4 without either file guessing.
\usepackage{xr}

%%  This manuscript quotes a lot of file paths in \texttt, and a typewriter
%%  path is one unbreakable word to TeX, so a paragraph that ends on one runs
%%  into the margin.  Two settings fix it without touching the prose: let a
%%  typewriter word hyphenate, and give TeX a last-resort stretch so it
%%  loosens a paragraph rather than overflowing it.  The first is what the
%%  `hyphenat' package's htt option does, written out here so that the twenty
%%  font-shape probes it performs at load time stay out of the build log --
%%  a clean log is the point of having one.
\makeatletter
\DeclareRobustCommand\texttt[1]{%
  {\ttfamily\hyphenchar\font=`\-\relax #1}}
\makeatother
\setlength{\emergencystretch}{3em}

%%  Twenty-five tables in sixty-five pages overflowed LaTeX's float queue: two
%%  of them --- the decision-curve table and the pilot's --- were deferred
%%  THIRTY pages past the text that reads them, which for a referee is the
%%  same as not having them.  The barrier keeps a table inside the section
%%  that discusses it; the fractions let a page carry more float and less
%%  compulsory text, so the barrier has somewhere to put them.
\usepackage[section]{placeins}
\renewcommand{\topfraction}{0.9}
\renewcommand{\bottomfraction}{0.8}
\renewcommand{\textfraction}{0.07}
\renewcommand{\floatpagefraction}{0.7}
\setcounter{topnumber}{3}
\setcounter{totalnumber}{5}

%% GENERATED by scripts/make_numbers.py -- do not edit.
%% Every numeric literal in the manuscript is one of these.
\newcommand{\Dabs}{+0.0899}
\newcommand{\DabsHi}{+0.1256}
\newcommand{\DabsLo}{+0.0634}
\newcommand{\Rref}{0.350}
\newcommand{\RrefKind}{bounded}
\newcommand{\VCorruptDelta}{-0.0357}
\newcommand{\VDuplicateDelta}{-0.0003}
\newcommand{\VPopCommon}{+0.1055}
\newcommand{\VPopRare}{+0.1039}
\newcommand{\VStaleDelta}{+0.0000}
\newcommand{\VTone}{+0.1835}
\newcommand{\VToneLo}{+0.1650}
\newcommand{\VTonehi}{+0.1909}
\newcommand{\VTtwoGroup}{+0.1033}
\newcommand{\VTtwoGroupHi}{+0.1106}
\newcommand{\VTtwoGroupLo}{+0.0762}
\newcommand{\VTtwoKnow}{-0.0032}
\newcommand{\VTtwoKnowHi}{+0.0013}
\newcommand{\VTtwoKnowLo}{-0.0079}
\newcommand{\agreeTolerance}{$1.1 \times 10^{-16}$}
\newcommand{\auditBJustifiedPct}{18\%}
\newcommand{\auditBStatedPct}{85\%}
\newcommand{\auditBoundHi}{48\%}
\newcommand{\auditBoundLo}{9\%}
\newcommand{\auditHalfWidthPct}{18 percentage points}
\newcommand{\auditKappa}{0.20}
\newcommand{\auditMJustifiedPct}{5\%}
\newcommand{\auditMaxWeight}{1.00}
\newcommand{\auditPrec}{0.481}
\newcommand{\auditPrecReg}{0.301}
\newcommand{\auditRangeBaselinePct}{34\%}
\newcommand{\auditRangeMetricPct}{15\%}
\newcommand{\auditRangePopulationPct}{0\%}
\newcommand{\auditRangeThresholdPct}{5\%}
\newcommand{\auditSens}{0.474}
\newcommand{\auditSensHi}{0.602}
\newcommand{\auditSensLo}{0.347}
\newcommand{\auditSensReg}{0.569}
\newcommand{\auditSpec}{0.911}
\newcommand{\auditSpecReg}{0.769}
\newcommand{\auditThetaStatedPct}{27\%}
\newcommand{\biasNoisyBefore}{-0.0353}
\newcommand{\biasOracleMax}{0.0405}
\newcommand{\cardF}{3{,}019}
\newcommand{\casePrevalence}{0.372}
\newcommand{\corruptLevel}{15\%}
\newcommand{\coverageNaiveMedian}{90.5\%}
\newcommand{\coverageNaiveMin}{5.5\%}
\newcommand{\coverageNestedBasicMedian}{92.0\%}
\newcommand{\coverageNestedBasicMin}{7.0\%}
\newcommand{\coverageNestedBcMedian}{91.5\%}
\newcommand{\coverageNestedMedian}{94.2\%}
\newcommand{\coverageNestedMin}{3.5\%}
\newcommand{\coverageNoisyAfter}{7.0\%}
\newcommand{\coverageNoisyBefore}{7.0\%}
\newcommand{\crossTargetAgreement}{45\%}
\newcommand{\doiBPICfourteen}{\url{https://doi.org/10.4121/uuid:c3e5d162-0cfd-4bb0-bd82-af5268819c35}}
\newcommand{\doiBPICthirteen}{\url{https://doi.org/10.4121/uuid:500573e6-accc-4b0c-9576-aa5468b10cee}}
\newcommand{\exCells}{480}
\newcommand{\exMisreportPct}{24.9\%}
\newcommand{\exRmax}{0.947}
\newcommand{\exRmin}{0.200}
\newcommand{\exRows}{30{,}162}
\newcommand{\exSbaselinePct}{25.9\%}
\newcommand{\exSmetricPct}{25.9\%}
\newcommand{\exSpopulationPct}{7.2\%}
\newcommand{\exampleLines}{121}
\newcommand{\flipLearnerPct}{31.3\%}
\newcommand{\flipMetricPct}{19.7\%}
\newcommand{\flipQualityPct}{20.4\%}
\newcommand{\flipRungPct}{22.2\%}
\newcommand{\flipSplitPct}{33.5\%}
\newcommand{\fvVersion}{0.2.0}
\newcommand{\gridHi}{0.80}
\newcommand{\gridLo}{0.05}
\newcommand{\kmBaseAuc}{0.648}
\newcommand{\kmOrphanPopulated}{100\%}
\newcommand{\lambdaHalf}{0.64}
\newcommand{\lambdaHalfPct}{64\%}
\newcommand{\limitNoisyAfter}{+0.1135}
\newcommand{\limitNoisyBefore}{+0.1135}
\newcommand{\maxCardBZero}{200}
\newcommand{\maxTMedian}{4.03}
\newcommand{\maxTMedianRatio}{2.06}
\newcommand{\meanNoisyEstimate}{+0.0781}
\newcommand{\minCardF}{50}
\newcommand{\minGPresent}{50\%}
\newcommand{\minPerStratum}{25}
\newcommand{\minReuseF}{two}
\newcommand{\minTrainCoverage}{50\%}
\newcommand{\misreportMaxPct}{57.7\%}
\newcommand{\misreportMedianPct}{35.3\%}
\newcommand{\misreportMinPct}{5.6\%}
\newcommand{\nAgreeDisagree}{0}
\newcommand{\nAgreeQuantities}{266}
\newcommand{\nAuditApplicable}{3}
\newcommand{\nAuditEffective}{29.6}
\newcommand{\nAuditEligible}{54.9}
\newcommand{\nAuditFrame}{600}
\newcommand{\nAuditFullText}{369}
\newcommand{\nAuditNoFullText}{231}
\newcommand{\nAuditOutSampled}{120}
\newcommand{\nAuditRequired}{196}
\newcommand{\nAuditSampled}{600}
\newcommand{\nAuditScreenedIn}{54}
\newcommand{\nAuditScreenedOut}{315}
\newcommand{\nAuditStrata}{77}
\newcommand{\nAuditTarget}{2{,}400}
\newcommand{\nAxes}{nine}
\newcommand{\nBandResolvedRaw}{2{,}046}
\newcommand{\nBandResolvedRecentred}{2{,}163}
\newcommand{\nCaseTest}{13{,}637}
\newcommand{\nCaseTrain}{31{,}818}
\newcommand{\nClosedBefore}{5}
\newcommand{\nCohort}{45{,}455}
\newcommand{\nCohortCorpus}{46{,}606}
\newcommand{\nCondBeneficial}{11}
\newcommand{\nCondHarmful}{1}
\newcommand{\nConditions}{8}
\newcommand{\nConfirmReject}{5}
\newcommand{\nCorpusLogs}{21}
\newcommand{\nDomains}{6}
\newcommand{\nDrawRows}{1{,}791{,}720}
\newcommand{\nFieller}{10{,}584}
\newcommand{\nFiellerUnbounded}{5{,}282}
\newcommand{\nFolds}{five}
\newcommand{\nGrid}{31}
\newcommand{\nHarmfulPointwise}{0}
\newcommand{\nHarmfulSimultaneous}{0}
\newcommand{\nLearners}{4}
\newcommand{\nLogs}{13}
\newcommand{\nLogsBothTargets}{6}
\newcommand{\nMeanOptimal}{19}
\newcommand{\nMetrics}{36}
\newcommand{\nMisreportOverQuarter}{13}
\newcommand{\nMisreportOverTen}{18}
\newcommand{\nOneNumberOptimal}{9}
\newcommand{\nOrphans}{94{,}250}
\newcommand{\nPairs}{19}
\newcommand{\nPairsBanded}{19}
\newcommand{\nPenaltyDecades}{five}
\newcommand{\nPropN}{1{,}800}
\newcommand{\nPropP}{200}
\newcommand{\nQuality}{10}
\newcommand{\nResultEstates}{18}
\newcommand{\nResultPairs}{12}
\newcommand{\nScaleReps}{\textbf{??}}
\newcommand{\nSeeds}{three}
\newcommand{\nSignChanging}{4}
\newcommand{\nSignVaries}{18}
\newcommand{\nSimReps}{200}
\newcommand{\nSimSparse}{200}
\newcommand{\nSizeGrid}{six}
\newcommand{\nSplits}{6}
\newcommand{\nSurfaceRows}{262{,}656}
\newcommand{\nTests}{108}
\newcommand{\nTrainOnly}{122}
\newcommand{\nTrainOnlyPassed}{122}
\newcommand{\nUniformlyBeneficial}{1}
\newcommand{\nUnresolvedPairs}{2}
\newcommand{\nWorkedBefore}{41{,}413}
\newcommand{\nWorlds}{six}
\newcommand{\noisyGapLarge}{\textbf{??}}
\newcommand{\noisyGapLoose}{\textbf{??}}
\newcommand{\noisyGapSmall}{\textbf{??}}
\newcommand{\noisyGapTight}{\textbf{??}}
\newcommand{\pThreeRawA}{0.421}
\newcommand{\pThreeRawB}{0.284}
\newcommand{\pThreeSqrtA}{0.232}
\newcommand{\pThreeSqrtB}{0.390}
\newcommand{\prevHi}{0.95}
\newcommand{\prevLo}{0.05}
\newcommand{\propOneApGap}{0.988}
\newcommand{\propOneAucGap}{exactly zero}
\newcommand{\propOneSweepAp}{0.998}
\newcommand{\propOneSweepAuc}{$4.4 \times 10^{-15}$}
\newcommand{\propTwoCancelShift}{machine precision}
\newcommand{\propTwoNonRankShift}{11.696}
\newcommand{\propTwoRankShift}{exactly zero}
\newcommand{\qClean}{+0.1673}
\newcommand{\qCommonHalf}{+0.1055}
\newcommand{\qCorruptThirty}{-0.0699}
\newcommand{\qDuplicateAll}{-0.0038}
\newcommand{\qLagHalfDelta}{-0.0031}
\newcommand{\qLagHalfPopulated}{94.1\%}
\newcommand{\qRandomHalf}{+0.0951}
\newcommand{\qRareHalf}{+0.1039}
\newcommand{\qStaleDelta}{+0.0000}
\newcommand{\qStaleKeepsDelta}{+0.0000}
\newcommand{\ratioBase}{2.3}
\newcommand{\ratioCheap}{9.0}
\newcommand{\ratioDear}{0.7}
\newcommand{\regretMajority}{0.0297}
\newcommand{\regretMajorityAdj}{0.0197}
\newcommand{\regretMajorityGainAdjPct}{18.3\%}
\newcommand{\regretMajorityGainPct}{21.9\%}
\newcommand{\regretMeanRule}{0.0242}
\newcommand{\regretOneNumber}{0.0380}
\newcommand{\regretOneNumberAdj}{0.0241}
\newcommand{\regretOneNumberExcessPct}{57.2\%}
\newcommand{\regretOracle}{0.0242}
\newcommand{\regretUniformRule}{0.0489}
\newcommand{\releaseTag}{v19.0}
\newcommand{\resultWitnessMax}{5}
\newcommand{\resultWitnessMin}{0}
\newcommand{\rhoMedian}{0.200}
\newcommand{\shiftMedian}{-0.0001}
\newcommand{\shiftNegativePct}{56\%}
\newcommand{\simNoiseRate}{20\%}
\newcommand{\simPrevalence}{0.02}
\newcommand{\sizeHi}{\textbf{??}}
\newcommand{\sizeLo}{\textbf{??}}
\newcommand{\sobolInteractionPct}{30.0\%}
\newcommand{\sobolLearnerPct}{2.6\%}
\newcommand{\sobolMaxAxis}{the metric}
\newcommand{\sobolMaxPct}{10.1\%}
\newcommand{\sobolMetricPct}{10.1\%}
\newcommand{\sobolQualityPct}{2.7\%}
\newcommand{\sobolRungPct}{1.9\%}
\newcommand{\sobolSplitPct}{9.4\%}
\newcommand{\sobolTargetPct}{2.1\%}
\newcommand{\sobolTopAxis}{the split}
\newcommand{\sobolTopPair}{BPIC19/duration}
\newcommand{\sobolTopPct}{50.7\%}
\newcommand{\sparseGapLarge}{\textbf{??}}
\newcommand{\sparseGapSmall}{\textbf{??}}
\newcommand{\spreadHeadline}{0.091 AUC at the median pair and 0.298 at the widest}
\newcommand{\spreadMax}{0.298}
\newcommand{\spreadMedian}{0.091}
\newcommand{\targetAgreement}{51\%}
\newcommand{\targetAgreementHi}{66\%}
\newcommand{\targetAgreementLo}{43\%}
\newcommand{\thetaBase}{0.30}
\newcommand{\thetaCheap}{0.10}
\newcommand{\thetaDear}{0.60}
\newcommand{\unseenMax}{87\%}
\newcommand{\unseenMin}{0\%}
\newcommand{\widthNoisyBefore}{0.0367}
\newcommand{\zenodoDOI}{the archived release cited in the data-availability statement (DOI reserved, inserted at proof)}


\externaldocument{supplement}

\newtheorem{proposition}{Proposition}
\newtheorem{lemma}{Lemma}
\newtheorem{corollary}{Corollary}
\newtheorem{definition}{Definition}
\newtheorem{remark}{Remark}
\theoremstyle{definition}
\newtheorem{assumption}{Assumption}

\journal{Information Systems}

\begin{document}

\begin{frontmatter}

\title{Specification Surfaces for Incremental Predictive Performance:
Estimation, Uncertainty, and Multi-Log Evaluation}

\author[ind]{Sudhir Dixit}
\ead{sudhir.dixit1@gmail.com}
\address[ind]{Independent Researcher}

\begin{abstract}
What a recorded field is worth to a prediction is reported as one number, but
that number is a functional of rarely stated choices: baseline, pipeline,
target, split, decision time, register quality, metric and operating point. We
define incremental predictive performance as a \emph{specification surface}
$V_s(f)$ over a declared design space under a declared measure, and give four
objects for reporting one: an exact functional-ANOVA decomposition of its
variance; simultaneous max-$t$ bands over the \emph{whole} surface a claim
ranges over, from a pipeline-refitting moving-block bootstrap; \emph{resolution
regions} with a robustness index; and \emph{specification regret}. On \nPairs\ log--target
pairs from \nLogs\ public event logs, over an audited
\nDeclaredAdmissibleScalar-cell space, the baseline alone moves the increment
by \spreadMedian\ AUC at the median pair, and a one-number report misstates
the sign of \misreportPooledResolvedPct\ of resolved cells ---
\misreportPooledItsmPct\ on the \nItsmPairs\ maintained registers. Of the
\nPairsResolvingAny\ pairs resolving any sign, \nPairsResolveBothSigns\ resolve
both. On \nModelsDcaExcluded\ of \nModelsDca\ models a threshold is not a
cost ratio at all. In a case study on \nCohort\ bank incidents, a register's
value turns on the decision time and its sign on the target.
\end{abstract}

%%  Six keywords: the Elsevier template caps them at six and an earlier
%%  version listed seven.
\begin{keyword}
specification surface \sep incremental predictive performance \sep
predictive process monitoring \sep simultaneous inference \sep
sensitivity analysis \sep configuration management database
\end{keyword}

\end{frontmatter}

%%  Highlights.  Elsevier collects these in a separate file; they are printed
%%  here as well so that a reviewer holding only the PDF can see them.  The
%%  text is generated into numbers.tex from submission/highlights.txt by
%%  scripts/make_numbers.py, and scripts/check_highlights.py enforces the
%%  three-to-five count and the eighty-five-character limit, so the two
%%  copies cannot drift and neither can be over length.
\section*{Highlights}
\begin{itemize}\itemsep0pt
\highlightsList{}
\end{itemize}

\section{Introduction}
\label{sec:intro}

A claim that a recorded field is worth something is a comparison. Its second
arm --- performance \emph{without} the field --- is a choice the analyst
makes, and so are the learner that makes the comparison, the encoding it uses,
the outcome it predicts, the way the data are split in time, whether the field
is available at the moment the decision is taken, how completely the register
behind the field is populated, the metric the difference is read in, and the
operating point at which that metric is evaluated. Each is routinely left
implicit, and the result is reported as one number.

This paper's position is not that such a number depends on those choices ---
that is known, and Section~\ref{sec:related} says by whom --- but that the
dependence is large enough, on real decisions, to change the sign of the
answer; that the right object to report is therefore a surface over a
\emph{declared} design space rather than a point; and that a surface admits
estimation and inference of its own.

\paragraph{The estimand}
Write a specification $s$ for a complete assignment to every design axis, and
\begin{equation}
V_s(f) \;=\; m_\theta\!\left(A_s(D_{\mathrm{train}}, B \cup \{f_q\}),\,
D_{\mathrm{test}}\right)
\;-\;
m_\theta\!\left(A_s(D_{\mathrm{train}}, B),\, D_{\mathrm{test}}\right),
\label{eq:estimand}
\end{equation}
where $A_s$ is the fitting procedure (learner, encoding, hyperparameters,
calibration), $B$ the baseline feature set, $f_q$ the field under a register
quality mechanism $q$, $m$ the metric and $\theta$ the operating point, and
where $D_{\mathrm{train}}, D_{\mathrm{test}}$ are fixed by the split rule and
the cohort. The \emph{specification surface} is the map
$s \mapsto V_s(f)$ over a design space $\mathcal{S}$ that the analyst declares
and the reader can enumerate; Figure~\ref{fig:speccurve} draws one.

We say \emph{incremental predictive performance} and not \emph{value}:
\eqref{eq:estimand} measures how much better a model predicts, not what the
field costs to acquire or maintain, whether acting on the prediction changes
the outcome, or what an organisation would pay for either.
Section~\ref{sec:dca} converts the increment into net benefit at a declared
exchange rate, which is as close to economic value as this evidence gets.

\begin{figure}[!t]
\centering
\includegraphics[width=\linewidth]{figS1_speccurve.png}
\caption{The specification curve for the configuration item on BPIC14, at the
registered handover target. \textbf{(a)} Each point is one admissible
specification --- a complete assignment to the pipeline, the baseline rung,
the register-quality condition, the split and the metric --- sorted by the
increment. \textbf{(b)} Which level of which axis each specification occupies.
The intercept-only rung is excluded. Every point is a point estimate; the
simultaneous bands that decide the region labels are in
Table~\ref{tab:calbands}. \textbf{The five scalar instruments are plotted on
one axis and their increments are not on a common scale}
(Section~\ref{sec:regretdef}); what the panel supports is the \emph{sign}
statement, which is unit-free, and not a comparison of magnitudes across
instruments. The panel shows the increment rather than the reduction, so a
cell where the register is worth almost nothing looks the same as one where
the free field absorbs almost everything.}
\label{fig:speccurve}
\end{figure}

\paragraph{Contributions}
\begin{enumerate}
\item \textbf{An estimand and four reporting objects} (Table~\ref{tab:objects}):
the surface over a declared design space, an exact functional-ANOVA
decomposition of its variance, simultaneous max-$t$ bands from a
pipeline-refitting bootstrap, \emph{resolution regions} with a robustness
index $\rho$, and \emph{specification regret}.
\item \textbf{What each is worth, measured rather than asserted.}
Section~\ref{sec:simband} measures the band's family-wise coverage against a
known answer and finds it short of nominal;
Section~\ref{sec:sca} reports where a specification curve's verdict and a
region label part company.
\item \textbf{A registered multi-log benchmark.} The surface is computed
through one frozen pipeline on every log--target pair a pre-registered
protocol admits --- \nPairs\ of them, not only the ones where the answer
resolves --- over an audited denominator, and every corpus statistic is
reported on all \nPairs\ and on the \nItsmPairs\ where the register is a
maintained one (Section~\ref{sec:family}).
\item \textbf{A register at three decision times.} Section~\ref{sec:cmdb}
shows that a decision time fixes which cases are in scope, what is known
about them and which fields are admissible at once, reconciles the two values
this log has been reported to give, and finds the difference is the
\emph{target}: over the cohort $\times$ target design the target carries
\shareTargetPct\ of the variance and the cohort \shareCohortPct.
\end{enumerate}

\begin{table}[t]
\centering
\footnotesize
\setlength{\tabcolsep}{4pt}%
\renewcommand{\arraystretch}{1.2}%
\begin{tabular}{>{\raggedright\arraybackslash}p{0.17\linewidth}>{\raggedright\arraybackslash}p{0.27\linewidth}cc>{\raggedright\arraybackslash}p{0.26\linewidth}}
\toprule
object & what it is for & \S & table & what it costs \\
\midrule
Sensitivity decomposition & how much of the disagreement each analysis choice explains, and how much belongs to no single one & \ref{sec:sobol} & \ref{tab:sobol} & exact; no refit. A M\"obius sum over the surface already computed \\
Simultaneous band & an interval that holds over every cell a claim quantifies over, not one cell at a time & \ref{sec:simbands} & \ref{tab:calbands} & the whole cost: one draw refits every arm of every cell \\
Resolution region and $\rho$ & which way the surface points where the data settle it, and how much of it they settle & \ref{sec:regionsdef} & \ref{tab:triple} & a function of the band; no further refit \\
Specification regret \emph{(a check, not an object)} & what choosing one specification costs against the best one, on a scale a decision can use --- it returns a null on this corpus & \ref{sec:regretdef} & \ref{tab:regret} & exact; no refit, on the surface already computed \\
\bottomrule
\end{tabular}
\caption{THE REPORTING OBJECTS. The three this paper supplies and the fourth it retains as a robustness check, what each is for, where it is defined, where it is reported on this corpus, and what it costs; the two middle columns are the section it is defined in and the table that reports it. Only the second costs model fits: a draw of the bootstrap refits every arm of every cell, which is why the surface an inference is taken over is smaller than the surface a point estimate is taken over (Section~\ref{sec:threesurfaces}). The other three are functions of a surface that has already been computed, so an analyst who has done the sweep has already paid for them.}
\label{tab:objects}
\end{table}


\paragraph{What this paper is not}
This is a retrospective study on public benchmarks, offered as an evaluation
methodology. It is not a procurement conclusion about configuration management
databases, not a systematic review of reporting practice, and not a report of
a deployed system; Section~\ref{sec:lim} says what each of those would have
required.

\paragraph{The supplementary document}
The protocol and its amendments, the proofs, the constructions of the four
reporting objects, the correction register, the results this article
summarises and the tables it does not print are in a supplementary document
whose sections, tables and figures carry an \textsf{S}. A reference here to
Section~S3 or Table~S12 is a reference to that document; the article is
self-contained without it.

\section{Related work}
\label{sec:related}

Every component used here is established. That a feature's apparent
importance depends on what else is in the model motivates conditional and
permutation importance \citep{strobl2008conditional,fisher2019models},
Shapley attributions \citep{lundberg2017unified,covert2020sage} and
nonparametric formulations of importance as a population quantity
\citep{williamson2021vimp}; the clinical-prediction literature has said for
two decades that the increment in AUC from adding a predictor depends on the
existing model \citep{pencina2008evaluating,cook2007roc}, is a poor guide to
clinical usefulness \citep{cook2007roc,vickers2008decision} and is routinely
misread \citep{kerr2011evaluating}. Decision-analytic evaluation supplies the
exchange rate \citep{vickers2006dca,vickers2016net,vickers2019simple}.
Multiverse and specification-curve analysis supply the sweep
\citep{steegen2016increasing,simonsohn2020specification,patel2015assessment},
underspecification and variance studies the observation that pipelines differ
\citep{damour2022underspecification,bouthillier2021variance}, and data
valuation a different question about a datum's worth
\citep{ghorbani2019data,jia2019towards}. On event logs, outcome-oriented
predictive monitoring supplies the task
\citep{teinemaa2019outcome,verenich2019survey,senderovich2017intercase},
encoding studies one axis \citep{liu2025encoding,carka2022effortaware}, and
data-quality work the vocabulary for a degraded register
\citep{suriadi2017imperfection,mohammed2025dataquality,comuzzi2025labels}.
Leakage discipline is standard
\citep{kaufman2012leakage,kapoor2023leakage,weytjens2022leakage}.
Supplement~\ref{app:related} sets out each strand and what this paper takes
from it.

\paragraph{The novelty gap, precisely}
Every component used here is established and none is claimed. What none of
them supplies is the object this paper reports, and the gap is in four
places. \emph{The estimand is not declared completely}: the axes are
discussed one at a time, and the decision time and the register's quality as
preliminaries rather than as arguments. \emph{The measure is left implicit}:
a specification curve is summarised uniformly over whatever cells were
enumerated, a choice that moves the answer (Section~\ref{sec:sobol}).
\emph{Inference is pointwise or absent}: a curve is a distribution of point
estimates and a permutation test on its median, which licenses no statement
about a sub-region. \emph{Nothing measures what the conventional report
costs}: the alternative to one number is offered as a display, not as a
decision rule with a loss. Of the four objects that close these, the
decomposition and the band are adaptations; \emph{resolution regions} with a
robustness index and \emph{specification regret} are, as far as we are aware,
new. Section~\ref{sec:sca} measures where a curve's verdict and a region
label part company rather than asserting a difference.

The contribution is deliberately \emph{not} ``people usually report one
number''. It is that a claim of this form is a functional of a design, that
the functional can be estimated and its uncertainty quantified, and that on
this corpus it is far from constant.


